\documentclass[10pt]{amsart}
\usepackage[a4paper,margin=23mm]{geometry}
\usepackage{amsmath,amssymb,amsthm,mathtools,hyperref,graphicx,xspace,letltxmacro,xpatch}
\usepackage[english]{babel}

\usepackage[scr=rsfs,cal=euler]{mathalfa}
\usepackage{mathtools}
\mathtoolsset{showonlyrefs}
\multlinegap0pt



\newcommand\mto{\mathchoice{\longmapsto}{\mapsto}{\mapsto}{\mapsto}}
\newcommand{\Psfrac}[2]{(\sfrac{#1}{#2})}

\usepackage{bbm}
\let\bm\mathbbm

\theoremstyle{plain}
\newtheorem{theo}{Theorem}[section]
\newtheorem{prop}[theo]{Proposition}
\newtheorem{coro}[theo]{Corollary}
\newtheorem{lemm}[theo]{Lemma}
\newtheorem{conj}[theo]{Conjecture}
\newtheorem{princ}[theo]{Principle}
\newtheorem*{ques}{Question}
\newtheorem{ex}[theo]{Example}
\theoremstyle{definition}
\newtheorem{def1}[theo]{Definition}
\theoremstyle{remark}
\newtheorem{rema}[theo]{Remark}

\newcommand{\Op}{\operatorname{Op}}
\newcommand{\U}{\mathcal{U}}
\newcommand{\nwc}{\newcommand}
%\nwc{\eps}{\epsilon}
%\nwc{\ep}{\epsilon}
%\nwc{\vareps}{\varepsilon}
%\nwc{\Oph}{\operatorname{Op}_\hbar}
%\nwc{\la}{\langle}
%\nwc{\ra}{\rangle}
%
%\nwc{\mf}{\mathbf} %Latex (as in \bf not tilted math letters)
%\nwc{\blds}{\boldsymbol} %Latex
\nwc{\ml}{\mathcal} %Latex
%
%\nwc{\defeq}{\stackrel{\rm{def}}{=}}

\nwc{\cE}{\ml{E}}
\nwc{\cN}{\ml{N}}
\nwc{\cO}{\ml{O}}
\nwc{\cP}{\ml{P}}
\nwc{\cU}{\ml{U}}
\nwc{\cV}{\ml{V}}
\nwc{\cW}{\ml{W}}
\nwc{\tU}{\widetilde{U}}
\nwc{\IN}{\mathbb{N}}
\nwc{\IR}{\mathbb{R}}
\nwc{\IZ}{\mathbb{Z}}
\nwc{\IC}{\mathbb{C}}
\nwc{\IT}{\mathbb{T}}
\nwc{\IS}{\mathbb{S}}
\nwc{\tP}{\widetilde{P}}
\nwc{\tPi}{\widetilde{\Pi}}
\nwc{\tV}{\widetilde{V}}
\nwc{\supp}{\operatorname{supp}}
\nwc{\rest}{\restriction}
\let \d \relax
\nwc{\d}{\partial}
\nwc{\Cor}{\mathscr{C}}

\nwc{\todo}[1]{$\clubsuit$ {\tt #1}}
\DeclareMathOperator{\vect}{span}
\DeclareMathOperator{\Vol}{Vol}
\DeclareMathOperator{\sgn}{sgn}
\DeclareMathOperator{\Hess}{Hess}
\DeclareMathOperator{\Res}{Res}
\let \div \relax \newcommand{\div}{\operatorname{div}}

\renewcommand{\Im}{\operatorname{Im}}
\renewcommand{\Re}{\operatorname{Re}}

\usepackage{xfrac,xcolor,graphicx}\graphicspath{{./figures/}}\providecommand{\skpt}{\par\smallskip\noindent}\providecommand{\pointir}{.\ }\providecommand{\eg}{e.g.}\providecommand{\ie}{i.e.}\providecommand{\resp}{resp.}\providecommand{\cf}{cf.}\providecommand{\etc}{etc.}\let\Subsection\subsection\let\Subsubsection\subsubsection
